% =====================================================================
\section{GaN: el transistor de alta movilidad (HEMT)}
\label{sec:gan}
% =====================================================================

\diapo{36} Con GaN se pueden hacer diodos parecidos a los de SiC, pero no es su interés principal.
Su dispositivo estrella es el \textbf{HEMT} (\emph{High Electron Mobility Transistor}), un
\textbf{transistor de heterounión}: funciona como un MESFET (un JFET con compuerta metálica, de
Schottky, en lugar de semiconductora), pero su canal no es una capa dopada sino un \textbf{gas
bidimensional de electrones} (2DEG) confinado en la interfaz entre dos semiconductores distintos. La
conductividad de ese canal es altísima y se logra con capacidades parásitas muy bajas.

\begin{ojo}[``Gas bidireccional'' \diapo{36}]
La diapo~36 dice ``gas \emph{bidireccional} de electrones''. Es gas \textbf{bidimensional} (2D gas,
2DEG): los electrones están confinados en un pozo de unos pocos nanómetros en la dirección
perpendicular a la interfaz y se mueven libremente en las otras dos. Las diapos~37 y 38 lo llaman
correctamente.
\end{ojo}

\subsection{La heterounión y el 2DEG}

Cuando se apoya un semiconductor de gap grande (AlGaN, o AlGaAs) sobre uno de gap menor (GaN, o
GaAs), la banda de conducción tiene un escalón $\Delta E_c$ en la interfaz. Del lado del material de
gap chico se forma un \textbf{pozo triangular} cuyo fondo queda \emph{por debajo} de $E_F$: allí se
acumulan electrones, que sólo pueden moverse paralelos a la interfaz.

\begin{figure}[H]
\centering
\begin{tikzpicture}[x=1cm,y=1cm,>=Stealth]
  % regiones
  \fill[gris!40] (-1.0,-4.2) rectangle (0,2.4);
  \fill[naranja!15] (0,-4.2) rectangle (2.2,2.4);
  \fill[azul!6] (2.2,-4.2) rectangle (9,2.4);
  \node[font=\small] at (-0.5,2.65) {metal};
  \node[font=\small] at (1.1,2.65) {AlGaN};
  \node[font=\small] at (5.6,2.65) {GaN (sin dopar)};
  % E_F
  \draw[fermi] (-1.0,0) -- (9,0) node[right,black,font=\small] {$E_F$};
  % Ec
  \draw[bandac] (0,1.9) -- (2.2,0.5) -- (2.2,-0.45) .. controls (2.6,-0.15) and (3.4,0.45) .. (5,0.65) -- (9,0.7);
  % Ev
  \draw[bandav] (0,-2.1) -- (2.2,-3.5) -- (2.2,-3.85) .. controls (2.6,-3.55) and (3.4,-2.95) .. (5,-2.75) -- (9,-2.7);
  % 2DEG
  \fill[azul!60,opacity=0.6] (2.2,-0.45) .. controls (2.35,-0.32) and (2.5,-0.2) .. (2.62,-0.1) -- (2.62,0) -- (2.2,0) -- cycle;
  \foreach \x/\y in {2.3/-0.15,2.42/-0.08,2.33/-0.32} \node[elec,minimum size=3.5pt] at (\x,\y) {};
  \draw[<-] (2.5,-0.12) -- (3.6,-0.9) node[right,font=\small,align=left] {2DEG: pozo triangular\\bajo $E_F$ ($\sim\SI{1e13}{cm^{-2}}$)};
  \draw[<->] (2.05,0.5) -- node[left,font=\small] {$\Delta E_c$} (2.05,-0.45);
  \node[azul,font=\small] at (7.5,1.0) {$E_c$};
  \node[rojo,font=\small] at (7.5,-2.4) {$E_v$};
  % cargas de polarización
  \foreach \y in {1.8,1.2,0.6,-1.8,-2.4,-3.0} \node[font=\scriptsize,verde!50!black] at (2.0,\y-0.15) {$+$};
  \node[font=\scriptsize,verde!50!black,align=center] at (1.0,-1.0) {carga de\\polarización\\$+\sigma_{\mathrm{pol}}$};
\end{tikzpicture}
\caption{Diagrama de bandas esquemático de una heterounión AlGaN/GaN bajo la compuerta. La
polarización del cristal deja una carga positiva fija en la interfaz; el campo resultante inclina la
banda del AlGaN y crea, del lado del GaN, un pozo triangular por debajo de $E_F$ donde se acumula el
2DEG, \textbf{sin necesidad de dopar} \diapo{38}.}
\label{fig:hemt-bandas}
\end{figure}

\paragraph{¿De dónde salen los electrones?}
\begin{itemize}
  \item \textbf{AlGaAs/GaAs} (HEMT tradicional, \diapo{37}): por \emph{dopaje modulado}. Se dopa el
  AlGaAs con donores; sus electrones caen al pozo del GaAs, que queda sin dopar. Los electrones quedan
  \emph{separados espacialmente} de los donores ionizados que los generaron: casi no hay dispersión
  por impurezas y la movilidad es muy alta (\SI{8500}{cm^2/V.s} a \SI{300}{K}, mucho más a baja
  temperatura). De ahí el nombre.
  \item \textbf{AlGaN/GaN}: el GaN y el AlGaN tienen estructura wurtzita, \emph{polar}: tienen
  polarización espontánea, y el AlGaN, al crecer tensionado sobre el GaN, además polarización
  piezoeléctrica. La discontinuidad de polarización deja en la interfaz una carga fija
  $\sigma_{\mathrm{pol}}/q\sim\SI{1e13}{cm^{-2}}$ que atrae una cantidad igual de electrones. Ambas
  capas pueden estar sin dopar \diapo{38}.
\end{itemize}

\begin{ejemplo}[Resistencia del canal 2DEG]
La resistencia de hoja de un canal de espesor despreciable con densidad superficial $n_s$ es
$R_\square=1/(qn_s\mu)$. Para el 2DEG de AlGaN/GaN, $n_s=\SI{1e13}{cm^{-2}}$ y
$\mu=\SI{2000}{cm^2/V.s}$:
\[
  R_\square=\frac{1}{\num{1.6e-19}\cdot\num{1e13}\cdot2000}\approx\SI{310}{\ohm}/\square .
\]
Un canal de inversión de un MOSFET de Si con la misma $n_s$ y $\mu\approx\SI{400}{cm^2/V.s}$ (la
movilidad superficial es mucho menor que la de volumen, \munSi) tendría \SI{1.6}{k\ohm}/$\square$: el
2DEG conduce cinco veces mejor con la misma carga, y por lo tanto con la misma capacidad de compuerta.
\end{ejemplo}

\subsection{El HEMT: estructura y control}

\begin{figure}[H]
\centering
\begin{tikzpicture}[x=1cm,y=1cm]
  \fill[rojo!25] (0,0) rectangle (10,0.8);
  \fill[gris!25] (0,0.8) rectangle (10,1.3);
  \fill[azul!8] (0,1.3) rectangle (10,3.0);
  \fill[naranja!40] (1.6,3.0) rectangle (8.4,3.35);
  \fill[azul!40,rounded corners=3pt] (0.3,2.4) rectangle (1.8,3.35);
  \fill[azul!40,rounded corners=3pt] (8.2,2.4) rectangle (9.7,3.35);
  \draw (0,0) rectangle (10,3.35);
  \draw[azul,line width=1.6pt,dashed] (1.8,2.95) -- (8.2,2.95);
  \fill[gris] (0.4,3.35) rectangle (1.7,3.75); \fill[gris] (8.3,3.35) rectangle (9.6,3.75);
  \fill[gris] (3.2,3.35) rectangle (4.2,3.75);
  \node[font=\small] at (1.05,4.0) {S}; \node[font=\small] at (3.7,4.0) {G}; \node[font=\small] at (8.95,4.0) {D};
  \node[font=\small] at (1.05,2.75) {$n^+$}; \node[font=\small] at (8.95,2.75) {$n^+$};
  \draw[<-] (7.4,3.2) -- (10.4,2.2) node[right,font=\small,align=left] {AlGaN\\($\sim\SI{20}{nm}$,\\sin dopar)};
  \node[font=\small] at (5,2.0) {GaN (sin dopar)};
  \node[font=\small] at (5,1.05) {capas de transición};
  \node[font=\small] at (5,0.4) {sustrato: Si, SiC o zafiro};
  \draw[<-] (5.6,2.92) -- (6.4,2.45) node[right,font=\small] {2DEG};
  \draw[<->] (4.2,4.45) -- node[above,font=\scriptsize] {$L_{GD}$: soporta la tensión} (8.3,4.45);
\end{tikzpicture}
\caption{HEMT de AlGaN/GaN \diapo{38}. Es un dispositivo \textbf{lateral}: la corriente circula por el
2DEG en el plano, y la tensión de bloqueo la soporta la distancia compuerta--drenaje $L_{GD}$ (y no el
espesor, como en los verticales). El GaN se crece sobre Si (barato, obleas grandes), SiC (mejor
disipación) o zafiro, con capas de transición que acomodan la diferencia de red cristalina.}
\label{fig:hemt-estructura}
\end{figure}

La compuerta (Schottky) controla el 2DEG como la de un MESFET \diapo{37}: con $V_{GS}=0$ el canal
existe y el transistor conduce; \textbf{polarizando la compuerta en inversa} ($V_{GS}<0$) se vacía el
pozo y se puede disminuir, e incluso eliminar, el 2DEG. El HEMT básico es por lo tanto un dispositivo
de \textbf{deplexión}, \emph{normalmente cerrado} (normally-on), con $V_{TH}\approx\SIrange{-3}{-5}{V}$.

\begin{intuicion}[RF vs.\ potencia \diapo{38}]
Con esta tecnología se construyeron primero amplificadores de radiofrecuencia (estaciones base,
radares), donde que el transistor sea normally-on no importa: trabaja polarizado en la región
lineal. En electrónica de potencia, en cambio, una llave que conduce cuando falta la tensión de
compuerta es inaceptable (el mismo problema del JFET de SiC, diapo~33). Hay que conseguir HEMT
\emph{normalmente abiertos}.
\end{intuicion}

\subsection{Cómo hacer un HEMT normalmente abierto}

\begin{figure}[H]
\centering
\begin{tikzpicture}[x=1cm,y=1cm]
  \foreach \dx/\tit in {0/{(a) Schottky (deplexión)},5.2/{(b) compuerta $p$-GaN},10.4/{(c) MIS con recess}} {
    \begin{scope}[xshift=\dx cm]
      \fill[azul!8] (0,0) rectangle (4.4,1.4);
      \node[font=\small] at (2.2,2.6) {\tit};
    \end{scope}
  }
  % (a)
  \fill[naranja!40] (0,1.4) rectangle (4.4,1.75);
  \fill[gris] (1.6,1.75) rectangle (2.8,2.15);
  \draw[azul,line width=1.5pt,dashed] (0,1.33) -- (4.4,1.33);
  \node[font=\scriptsize] at (2.2,0.6) {2DEG continuo: conduce con $V_{GS}=0$};
  % (b)
  \begin{scope}[xshift=5.2cm]
    \fill[naranja!40] (0,1.4) rectangle (4.4,1.75);
    \fill[rojo!40] (1.6,1.75) rectangle (2.8,2.05);
    \fill[gris] (1.7,2.05) rectangle (2.7,2.3);
    \draw[azul,line width=1.5pt,dashed] (0,1.33) -- (1.5,1.33);
    \draw[azul,line width=1.5pt,dashed] (2.9,1.33) -- (4.4,1.33);
    \node[font=\scriptsize,align=center] at (2.2,0.6) {la $p$-GaN levanta la banda:\\sin 2DEG bajo la compuerta};
  \end{scope}
  % (c)
  \begin{scope}[xshift=10.4cm]
    \fill[naranja!40] (0,1.4) rectangle (1.5,1.75); \fill[naranja!40] (2.9,1.4) rectangle (4.4,1.75);
    \fill[gris!40] (1.4,1.25) rectangle (3.0,1.4); \fill[gris!40] (1.4,1.4) rectangle (1.5,1.85); \fill[gris!40] (2.9,1.4) rectangle (3.0,1.85);
    \fill[gris] (1.5,1.4) rectangle (2.9,1.95);
    \draw[azul,line width=1.5pt,dashed] (0,1.33) -- (1.4,1.33);
    \draw[azul,line width=1.5pt,dashed] (3.0,1.33) -- (4.4,1.33);
    \node[font=\scriptsize,align=center] at (2.2,0.6) {se elimina el AlGaN bajo la compuerta;\\óxido $+$ canal tipo MOS};
  \end{scope}
\end{tikzpicture}
\caption{Compuertas de HEMT de GaN (sólo la zona de compuerta; naranja: AlGaN, rojo: $p$-GaN, gris
claro: dieléctrico). (a) Compuerta Schottky: normally-on. (b) Compuerta de GaN tipo $p$: su
potencial interno vacía el 2DEG bajo la compuerta a $V_{GS}=0$; $V_{TH}\approx\SIrange{1}{2}{V}$.
(c) Compuerta MIS con el AlGaN removido: se forma un canal por acumulación como en un MOSFET.}
\label{fig:hemt-compuertas}
\end{figure}

La diapo~39 clasifica los HEMT de AlGaN/GaN comerciales por tipo de compuerta y comportamiento:

\begin{table}[H]
\centering\small
\begin{tabularx}{\textwidth}{>{\raggedright\arraybackslash}p{3.0cm}>{\raggedright\arraybackslash}p{2.6cm}>{\raggedright\arraybackslash}X>{\raggedright\arraybackslash}p{3.4cm}}
\toprule
Comportamiento & Compuerta & Uso & Fabricantes (diapo 39) \\
\midrule
normalmente cerrado & metálica (MES) & RF (amplificación) & Cree, Nitronex, TriQuint \\
normalmente abierto & metálica sobre $p$-GaN & electrónica de potencia (directo) & EPC, Panasonic, GaN Systems \\
normalmente cerrado & MOS (MIS) & electrónica de potencia \textbf{en cascodo} con MOSFET de Si (fig.~\ref{fig:cascodo}) & Transphorm, MicroGaN, IR \\
normalmente abierto & MOS (MIS) & en desarrollo (``¿?'' en la diapo) & --- \\
\bottomrule
\end{tabularx}
\caption{Tipos de HEMT de AlGaN/GaN para potencia \diapo{39}.}
\label{tab:hemt-tipos}
\end{table}

\paragraph{Ejemplo de Infineon \diapo{40}.} El dispositivo de la diapo (desarrollado por International
Rectifier, hoy parte de Infineon) es un HEMT \emph{de deplexión} con compuerta MIS: un dieléctrico bajo
la compuerta (``Metal Insulator Semiconductor'') que reduce mucho la fuga de compuerta respecto de una
Schottky, fabricado con procesos CMOS estándar sobre sustrato de Si. Es \textbf{normalmente cerrado} y
está pensado para usarse \textbf{en cascodo} con un MOSFET de Si de baja tensión, exactamente como el
JFET de SiC (figura~\ref{fig:cascodo}): el conjunto es una llave normally-off. (La línea CoolGaN actual
de Infineon usa compuerta $p$-GaN, normally-off.)

\subsection{Lo que ganan y lo que cuesta el GaN}

\begin{itemize}
  \item \textbf{Velocidad}: $\Qg$ y capacidades muy chicas $\Rightarrow$ conmutación en pocos
  nanosegundos, convertidores de \SIrange{1}{10}{MHz} con inductores y capacitores diminutos (los
  cargadores USB-C de alta potencia son de GaN).
  \item \textbf{Sin recuperación inversa}: no hay juntura $pn$ ni diodo de cuerpo; el HEMT conduce en
  inversa por el propio canal (con una caída de unos volts) y no almacena carga.
  \item \textbf{Limitaciones}: es lateral (hasta \SI{650}{V} típicos; los de \SI{1200}{V} recién
  aparecen), la compuerta $p$-GaN tolera poca tensión (\SIrange{6}{7}{V}) y exige drivers cuidados,
  y las trampas del material producen \emph{colapso de corriente} ($\RDSon$ dinámica: la resistencia
  sube justo después de bloquear alta tensión). Su conductividad térmica (\SI{1.3}{W/cm.K}) es menor
  que la del SiC, por lo que para alta potencia se lo crece sobre SiC.
\end{itemize}
